\newpage
\singlespacing
\begin{center}
	{\large Cultural Consumption, Culture-Based Talk, and Stratification}
\end{center}

\begin{abstract}
\small I analyze survey data from a probability sample of Americans to  for the first time document the predictors of what I call ``culture-based talk'': essentially, the use of cultural content to generate conversation with friends and acquaintances.  I argue that this phenomenon---central to many \emph{theories} of the connection between culture and stratification---deserves more \emph{empirical} attention. I show that rates of culture-based talk appear to be stratified by the markers of status (such as levels of schooling), and for some cultural domains (such as music), this stratification is independent of frequency of cultural engagement.  I close by outlining the theoretical and empirical implications of taking the phenomenon of culture-based talk seriously. 
\end{abstract}